\documentclass[11pt]{article}

\usepackage[margin=1in]{geometry}
\usepackage{setspace}
\usepackage{enumitem}
\usepackage{booktabs}
\usepackage{hyperref}
\usepackage[backend=biber, style=numeric]{biblatex}
\addbibresource{main.bib}

\hypersetup{colorlinks=true, linkcolor=blue, urlcolor=blue}
\onehalfspacing

\title{\textbf{Project Proposal}\\[0.5em]
\large Stabilizing Cortical Wetware via Artificial Immune System Algorithms  }
\author{Eliza Kishan \\
Department of Computer Science and Department of Biology \\
261211296 \\
COMP 599 - Natural Computation}
\date{\today}

\begin{document}

\maketitle

\begin{abstract}
Wetware is a new frontier in data processing, utilizing human neurons instead of basic silicon as the main data-processing
system. New data centers, such as the Biological Computing Facility in Singapore, promise new methods to approaching modern computational 
and biological problems. 
\end{abstract}

\section{Introduction}
Introduce the project topic and explain its broader context and significance.

\section{Problem Statement}
Clearly define the problem or question this project will address. Identify the
intended users, stakeholders, or affected communities.

\section{Objectives}
List the specific, measurable goals of the project.
\begin{enumerate}[label=\arabic*.]
	\item Objective one.
	\item Objective two.
	\item Objective three.
\end{enumerate}

\section{Background and Related Work}
Summarize relevant research, existing solutions, and prior work. Explain how this
project differs from or builds upon them, citing sources as appropriate.

\section{Proposed Approach}
Describe the planned design, methodology, algorithms, tools, and technologies.
Include diagrams, models, or equations if they clarify the approach.

\section{Scope and Limitations}
Define what is included in the project and identify assumptions, constraints, and
known limitations.

\section{Deliverables}
Describe the tangible outcomes that will be produced, such as software, datasets,
documentation, experiments, or a final report.

\section{Evaluation Plan}
Explain how success will be measured. Specify evaluation criteria, baselines,
datasets, experiments, user studies, or other validation methods.

\section{Project Schedule}
\begin{table}[ht]
	\centering
	\begin{tabular}{p{0.25\textwidth} p{0.6\textwidth}}
		\toprule
		\textbf{Timeframe} & \textbf{Milestone or Activity} \\
		\midrule
		Week 1--2 & Topic refinement and background research \\
		Week 3--4 & Design and implementation planning \\
		Week 5--8 & Development and experimentation \\
		Week 9--10 & Evaluation and revisions \\
		Week 11--12 & Final report and presentation \\
		\bottomrule
	\end{tabular}
	\caption{Proposed project schedule.}
\end{table}

\section{Resources and Budget}
Identify required software, hardware, data, personnel, and estimated costs.

\section{Risks and Mitigation}
\begin{table}[ht]
	\centering
	\begin{tabular}{p{0.35\textwidth} p{0.5\textwidth}}
		\toprule
		\textbf{Risk} & \textbf{Mitigation Strategy} \\
		\midrule
		Risk one & Planned response \\
		Risk two & Planned response \\
		\bottomrule
	\end{tabular}
	\caption{Potential risks and mitigation strategies.}
\end{table}

\section{Expected Contributions}
State the anticipated academic, technical, practical, or social contributions of
the project.

\section{Conclusion}
Restate the problem, proposed approach, expected outcomes, and project significance.

\begin{thebibliography}{9}
\bibitem{example}
Author, A. A. (Year). \textit{Title of source}. Publisher or Journal, volume(issue), pages.
\end{thebibliography}

\end{document}
